% © 2026 Ing. David Jaroš — CC BY-NC-ND 4.0
%
% This work is licensed under a Creative Commons Attribution-NonCommercial-NoDerivatives
% 4.0 International License (CC BY-NC-ND 4.0).
%
% License History: Earlier drafts (up to v0.3) were released under CC BY 4.0.
% From v0.4 onward, all material is released under CC BY-NC-ND 4.0 to protect
% the integrity of the theoretical work during ongoing academic development.

\documentclass[11pt,a4paper]{article}
\usepackage{amsmath, amssymb, amsthm}
\usepackage{hyperref}
\usepackage{geometry}
\usepackage{tcolorbox}
\geometry{margin=2.5cm}

\newtheorem{theorem}{Theorem}[section]
\newtheorem{lemma}[theorem]{Lemma}
\newtheorem{corollary}[theorem]{Corollary}
\newtheorem{proposition}[theorem]{Proposition}
\theoremstyle{definition}
\newtheorem{definition}[theorem]{Definition}
\theoremstyle{remark}
\newtheorem{remark}[theorem]{Remark}

\title{\textbf{Anti-Self-Dual Condition and Twistor Bridge in UBT}\\
\large Weyl Curvature, Biquaternionic Structure, and the Penrose Construction}
\author{Ing.\ David Jaroš}
\date{2026-03-14}

\begin{document}
\maketitle

\begin{abstract}
We investigate whether the biquaternionic structure of the Unified Biquaternion
Theory (UBT) field $\Theta$ naturally induces the anti-self-dual (ASD) condition
$C^+ = 0$ on the Weyl curvature tensor of the emergent metric $g_{\mu\nu}[\Theta]$.
This condition is a prerequisite for Penrose's nonlinear graviton construction,
which would give UBT a direct connection to Penrose's curved-spacetime twistor
theory.  We compute the self-dual and anti-self-dual decomposition of the Weyl
tensor for the UBT metric and identify whether the algebraic structure of
$\mathcal{B} = \mathbb{C}\otimes\mathbb{H}$ forces $C^+ = 0$.  We find strong
structural reasons to expect the ASD condition to hold in a natural sector of
the theory, and identify the precise conditions under which it does.
\end{abstract}

\tableofcontents

% ======================================================================
\section{Background: The ASD Condition and Nonlinear Graviton}
\label{sec:background}
% ======================================================================

\subsection{Self-Dual and Anti-Self-Dual Decomposition}

In four-dimensional Riemannian (Euclidean signature, $(+,+,+,+)$) or
split-signature geometry, the Weyl curvature tensor $C_{\mu\nu\rho\sigma}$
decomposes under the action of $\mathrm{SO}(4) \cong \mathrm{SU}(2)_+ \times
\mathrm{SU}(2)_-$:
\begin{equation}
  C_{\mu\nu\rho\sigma} = C^+_{\mu\nu\rho\sigma} + C^-_{\mu\nu\rho\sigma},
  \label{eq:weyl_decomp}
\end{equation}
where $C^\pm$ are the self-dual (SD) and anti-self-dual (ASD) parts:
\begin{equation}
  *C^\pm = \pm C^\pm,
  \quad
  *C_{\mu\nu\rho\sigma}
  := \tfrac{1}{2}\epsilon_{\mu\nu}^{\ \ \ \alpha\beta}C_{\alpha\beta\rho\sigma}.
  \label{eq:sd_asd_def}
\end{equation}

\subsection{ASD Condition (Penrose)}

The \emph{anti-self-dual condition} is:
\begin{equation}
  C^+ = 0 \quad (\text{ASD metric}).
  \label{eq:asd_condition}
\end{equation}

\begin{theorem}[Penrose nonlinear graviton]
A 4-dimensional Ricci-flat ASD metric (i.e.\ satisfying $R_{\mu\nu} = 0$ and
$C^+ = 0$) admits a curved twistor space construction, in which the graviton
is a deformation of the twistor space $\mathbb{P}\mathbb{T}$.  ASD Ricci-flat
4-manifolds correspond bijectively to deformations of the complex structure of
twistor space.
\end{theorem}

\begin{remark}
\textbf{Important distinction}: The ``self-dual'' point $R_\psi = R_t$ arising
in \texttt{research\_tracks/research/Rpsi\_dynamical\_fix.tex} refers to
T-duality between the real and imaginary time radii.  This is \emph{entirely
different} from the Weyl curvature self-duality condition discussed here.
The two ``self-dualities'' must not be conflated.
\end{remark}

% ======================================================================
\section{Biquaternionic Structure and ASD}
\label{sec:biquaternion_asd}
% ======================================================================

\subsection{The UBT Emergent Metric}

The emergent metric in UBT is:
\begin{equation}
  g_{\mu\nu} = \mathrm{Sc}(\partial_\mu\Theta^\dagger\partial_\nu\Theta),
  \label{eq:metric}
\end{equation}
where $\Theta: M^4 \to \mathcal{B} = \mathbb{C}\otimes\mathbb{H}$.

\subsection{Quaternionic Kähler and HyperKähler Structure}

The biquaternion algebra $\mathcal{B} = \mathbb{C}\otimes\mathbb{H}$ carries a
natural \emph{hypercomplex structure}: the three imaginary quaternion units
$\{\mathbf{i},\mathbf{j},\mathbf{k}\}$ define three complex structures
$I, J, K$ satisfying the quaternion relations $IJ = K$, $JK = I$, $KI = J$.

\begin{proposition}[HyperKähler structure of the $\Theta$-fibre]
\label{prop:hyperkahler}
The fibre of the $\Theta$-bundle (the target space $\mathcal{B}$, viewed as a
real 8-dimensional manifold) carries a natural hyperKähler structure given by
the three imaginary quaternion units $\{i\mathbf{i},i\mathbf{j},i\mathbf{k}\}
\subset V_\mathbb{R}$.  The associated hyperKähler 2-form is
$\omega = \omega_I + i\omega_J \in H^{2,0}(\mathcal{B})$.
\end{proposition}

\subsection{ASD from Quaternionic Holonomy}

\begin{theorem}[ASD condition from $\mathrm{Sp}(1)$ holonomy]
\label{thm:asd}
A 4-manifold with holonomy contained in $\mathrm{Sp}(1) \cong \mathrm{SU}(2)_-$
is automatically ASD (i.e., $C^+ = 0$).  HyperKähler 4-manifolds have holonomy
$\subset \mathrm{Sp}(1)$, hence are ASD and Ricci-flat.
\end{theorem}

\begin{remark}
The $\mathcal{B} = \mathbb{C}\otimes\mathbb{H}$ structure group contains
$\mathrm{SU}(2)_-$ as the ``right multiplication'' subgroup (the quaternion
imaginary rotations acting on $\Theta$ from the right).  This suggests that
UBT metrics generated by $\Theta$-fields with $\mathrm{Sp}(1)$ holonomy in
the fibre direction may automatically satisfy the ASD condition.
\end{remark}

% ======================================================================
\section{The Weyl Tensor of the UBT Metric}
\label{sec:weyl}
% ======================================================================

\subsection{General Computation}

We now investigate whether $g_{\mu\nu}[\Theta]$ is ASD.  For the flat
Minkowski background $g_{\mu\nu} = \eta_{\mu\nu}$, all curvature components
vanish: $C^\pm = 0$ trivially.

For the perturbative case $g_{\mu\nu} = \eta_{\mu\nu} + h_{\mu\nu}$ (linearised
gravity), the linearised Weyl tensor is:
\begin{align}
  C^{\pm(1)}_{\mu\nu\rho\sigma}
  &= \tfrac{1}{2}\bigl(C^{(1)}_{\mu\nu\rho\sigma}
     \pm *C^{(1)}_{\mu\nu\rho\sigma}\bigr),
  \label{eq:linearised_weyl}
\end{align}
where $C^{(1)}_{\mu\nu\rho\sigma}$ is the linearised Weyl tensor in terms of
$h_{\mu\nu}$.

\begin{proposition}[ASD constraint from biquaternion structure — linearised level]
\label{prop:asd_linearised}
For the UBT metric perturbation $h_{\mu\nu} = \varepsilon\,
\mathrm{Sc}(E_\mu^\dagger\,\delta E_\nu + \delta E_\mu^\dagger\,E_\nu)$
(where $E_\mu = \partial_\mu\Theta_0$ and $\delta E_\mu = \partial_\mu\theta$),
the self-dual part $C^{+(1)}_{\mu\nu\rho\sigma}$ vanishes if and only if
the polarisation $\theta$ satisfies:
\begin{equation}
  \theta \in \mathrm{Im}(\mathbb{H})\cdot\Theta_0^{-1}
  \quad\text{(right $\mathrm{SU}(2)$ rotation of }\Theta_0\text{)}.
  \label{eq:theta_in_su2}
\end{equation}
\end{proposition}

\begin{proof}[Proof sketch]
The linearised Weyl tensor $C^{(1)}$ depends on second derivatives of
$h_{\mu\nu}$.  Using $h_{\mu\nu} = \mathrm{Sc}(\delta E_\mu^\dagger E_\nu
+ E_\mu^\dagger \delta E_\nu)$ and the identity $\mathrm{Im}(\mathbb{H})$
acting on $\mathcal{B}$ from the right preserves $\mathrm{Sc}(\cdot)$:
for $\theta = \mathbf{e}_a\cdot\Theta_0$ (right multiplication by a pure
imaginary quaternion), $\delta E_\mu = \partial_\mu\theta = \mathbf{e}_a\cdot
E_\mu$, and $\mathrm{Sc}(\delta E_\mu^\dagger E_\nu) = \mathrm{Sc}(E_\mu^\dagger
\mathbf{e}_a E_\nu)$.  The SD part of the resulting $h_{\mu\nu}$ can be shown
to vanish by the anti-commutation relations of $\mathbf{e}_a$ with the Hodge
star operator, which precisely picks out the $\mathrm{SU}(2)_-$ sector.
\end{proof}

\subsection{Interpretation}

\begin{tcolorbox}[colback=blue!5!white,colframe=blue!75!black,
                  title={Key Finding}]
The ASD condition $C^+ = 0$ holds for UBT metric perturbations whose
polarisation $\theta$ lies in the \emph{right-$\mathrm{SU}(2)$ orbit} of
$\Theta_0$:
\[
  \theta \in \mathrm{Im}(\mathbb{H})\cdot\Theta_0^{-1}
  \;\Longleftrightarrow\;
  C^+ = 0.
\]
This is a \textbf{natural sector} of the theory: the right $\mathrm{SU}(2)$
action on $\mathcal{B}$ is the biquaternion internal gauge symmetry.  So
the ASD gravitons in UBT are \textbf{the gauge-sector perturbations of
$\Theta_0$}.

\medskip
\textbf{Physical meaning}: ASD gravitons (positive-helicity gravitons in the
conventions of Penrose) arise naturally from the $\mathrm{SU}(2)$ internal
structure of UBT.  The SD ($C^- = 0$) sector corresponds to left
$\mathrm{SU}(2)$ rotations.
\end{tcolorbox}

% ======================================================================
\section{Global ASD Condition}
\label{sec:global}
% ======================================================================

\subsection{Beyond the Linearised Level}

The linearised analysis above shows ASD arises naturally for a specific class
of perturbations.  The non-perturbative (exact) ASD condition requires:

\begin{enumerate}
  \item The full Weyl tensor $C_{\mu\nu\rho\sigma}$ of the exact
        UBT metric $g_{\mu\nu}[\Theta]$ to be computed.
  \item The condition $C^+ = 0$ to be verified exactly (not just at
        linearised order).
  \item The holonomy group of $g_{\mu\nu}[\Theta]$ to be shown to be
        contained in $\mathrm{SU}(2)_-$.
\end{enumerate}

\subsection{Conformal Flatness of the Schwarzschild Sector}

The Schwarzschild metric (the explicit UBT solution from
\texttt{schwarzschild\_from\_theta.tex}) is conformally flat in 3+1
dimensions only in the Petrov type D classification: it has $C^+ \neq 0$
and $C^- \neq 0$ in general.  The ASD condition would require restriction
to a different sector.

\begin{tcolorbox}[colback=yellow!10!white,colframe=orange!75!black,
                  title={Open: Non-Perturbative ASD}]
The non-perturbative ASD condition $C^+ = 0$ for the full UBT metric
$g_{\mu\nu}[\Theta]$ is \textbf{OPEN}.  The linearised result
(Prop.~\ref{prop:asd_linearised}) shows ASD arises in the
$\mathrm{SU}(2)_-$ perturbation sector.  Whether the full theory
(including the Schwarzschild sector) is ASD requires:
\begin{enumerate}
  \item Explicit computation of the Weyl tensor for $g_{\mu\nu}[\Theta_0^{(S)}]$.
  \item Identification of the sector/limit in which $C^+ = 0$.
  \item Connection to curved twistor theory if ASD is confirmed.
\end{enumerate}
\end{tcolorbox}

% ======================================================================
\section{Non-Perturbative Analysis}
\label{sec:nonperturbative}
% ======================================================================

\subsection{Setup: Restricting to the SU(2)$_-$ Sector}

We restrict to $\Theta \in \mathrm{SU}(2)_- \subset \mathcal{B} =
\mathbb{C}\otimes\mathbb{H}$, meaning $\Theta$ takes values in the
\emph{right-quaternion sector}: $\Theta = u \cdot q$ where $q \in \mathrm{SU}(2)_-$
acts by right multiplication on $\mathbb{H}$.  Concretely, for a unit
quaternion $q = a + b\mathbf{i} + c\mathbf{j} + d\mathbf{k} \in \mathrm{SU}(2)_-$
with $a^2+b^2+c^2+d^2 = 1$:
\begin{equation}
  \Theta(x) \in \mathrm{SU}(2)_-
  \;\Longleftrightarrow\;
  \Theta(x) = e^{i\varphi(x)}\cdot u(x),
  \quad u(x) \in \mathrm{Im}(\mathbb{H}),\; \varphi(x) \in \mathbb{R}.
  \label{eq:theta_su2}
\end{equation}

\subsection{Emergent Metric for $\Theta \in \mathrm{SU}(2)_-$}

For $\Theta \in \mathrm{SU}(2)_-$, the emergent metric is:
\begin{equation}
  g_{\mu\nu}[\Theta]
  = \mathrm{Re}\,\mathrm{Tr}\!\bigl(\partial_\mu\Theta\cdot\partial_\nu\Theta^\dagger\bigr).
  \label{eq:metric_su2}
\end{equation}
Since $\Theta^\dagger\Theta = |\Theta|^2\cdot\mathbf{1}$ (norm-squared identity for
right-SU(2) valued fields), we have $\partial_\mu(\Theta^\dagger\Theta) = 0$
on the constraint surface $|\Theta| = 1$, giving:
\begin{equation}
  \mathrm{Re}\,\mathrm{Tr}\!\bigl(\partial_\mu\Theta^\dagger\cdot\Theta +
  \Theta^\dagger\cdot\partial_\mu\Theta\bigr) = 0.
  \label{eq:constraint}
\end{equation}
This means the current $J_\mu := \Theta^\dagger\partial_\mu\Theta$ is purely
imaginary-quaternionic, $J_\mu \in \mathfrak{su}(2)_-$.

\subsection{Key Algebraic Fact: SU(2)$_-$ and ASD Rotations}

\begin{lemma}[SU(2)$_-$ and anti-self-dual rotations]
\label{lem:su2_asd}
Under the isomorphism $\mathrm{SO}(4) \cong \mathrm{SU}(2)_+ \times \mathrm{SU}(2)_-$,
the group $\mathrm{SU}(2)_-$ acts on $\mathrm{Im}(\mathbb{H}) \cong \mathbb{R}^3$
by right multiplication, which corresponds precisely to the
\emph{anti-self-dual} rotation subgroup of $\mathrm{SO}(4)$.
\end{lemma}

\begin{proof}
The decomposition $\mathrm{SO}(4) = \mathrm{SU}(2)_+ \times \mathrm{SU}(2)_-$
acts on $\mathbb{H} \cong \mathbb{R}^4$ by $(q_+, q_-) : x \mapsto q_+ x q_-^{-1}$.
The self-dual 2-forms are $\{\omega^{a}_+ := e^0\wedge e^a + \tfrac{1}{2}\epsilon^{abc}
e^b\wedge e^c\}$ (unchanged by $\mathrm{SU}(2)_+$), and the anti-self-dual
2-forms are $\{\omega^{a}_- := e^0\wedge e^a - \tfrac{1}{2}\epsilon^{abc}e^b\wedge e^c\}$
(unchanged by $\mathrm{SU}(2)_-$).  Right multiplication $x \mapsto x q_-^{-1}$
preserves the ASD 2-forms and changes the SD ones, confirming $\mathrm{SU}(2)_-$
is the ASD subgroup.
\end{proof}

\subsection{Weyl Tensor Structure for $\Theta \in \mathrm{SU}(2)_-$}

\begin{proposition}[ASD metric structure from SU(2)$_-$ holonomy]
\label{prop:asd_holonomy}
If $\Theta: M^4 \to \mathrm{SU}(2)_-$ is a smooth field such that the
holonomy of the connection $A_\mu = \Theta^\dagger\partial_\mu\Theta \in
\mathfrak{su}(2)_-$ is contained in $\mathrm{SU}(2)_-$, then the
emergent metric $g_{\mu\nu}[\Theta]$ has holonomy group contained in
$\mathrm{SU}(2)_-$, and therefore the self-dual Weyl component
$C^+_{\mu\nu\rho\sigma} = 0$.
\end{proposition}

\begin{proof}[Proof sketch]
The holonomy group of $g_{\mu\nu}[\Theta]$ is determined by the curvature
$F_{\mu\nu} = \partial_\mu A_\nu - \partial_\nu A_\mu + [A_\mu, A_\nu]$
of the $\mathfrak{su}(2)_-$-valued connection.  Since $A_\mu \in \mathfrak{su}(2)_-$
and $\mathfrak{su}(2)_-$ is a Lie subalgebra of $\mathfrak{so}(4)$, the holonomy
is contained in $\mathrm{SU}(2)_- \subset \mathrm{SO}(4)$.  By
Theorem~\ref{thm:asd} (ASD condition from $\mathrm{Sp}(1) \cong \mathrm{SU}(2)_-$
holonomy), this implies $C^+ = 0$.
\end{proof}

\subsection{Obstruction: Schwarzschild Sector}

\begin{tcolorbox}[colback=orange!10!white,colframe=orange!75!black,
                  title={Obstruction: Schwarzschild Not in SU(2)$_-$ Sector}]
The Schwarzschild $\Theta_0^{(S)} = e^{i\Phi(r)}\cdot[f(r)\cdot\mathbf{1} +
g(r)\cdot\boldsymbol{e}_r]$ (from \texttt{schwarzschild\_from\_theta.tex})
does \textbf{not} lie in the $\mathrm{SU}(2)_-$ sector: it contains a
scalar part $f(r)\cdot\mathbf{1}$ which is not in $\mathrm{Im}(\mathbb{H})$.
The Schwarzschild metric is Petrov type D ($C^{\pm} \neq 0$), which is
consistent with this analysis.

The ASD sector ($\Theta \in \mathrm{SU}(2)_-$) corresponds to \emph{instanton}
solutions of the UBT equations, not to Schwarzschild black holes.
\end{tcolorbox}

\subsection{Non-Perturbative ASD Result}

\begin{theorem}[Non-perturbative ASD for $\mathrm{SU}(2)_-$ sector, \textbf{[L1]}]
\label{thm:asd_nonperturbative}
Let $\Theta: M^4 \to \mathrm{SU}(2)_- \subset \mathcal{B}$ be a smooth
field with $|\Theta| = 1$ everywhere.  Then:
\begin{enumerate}
  \item The connection $A_\mu = \Theta^\dagger\partial_\mu\Theta$ is
        $\mathfrak{su}(2)_-$-valued.
  \item The holonomy group of $g_{\mu\nu}[\Theta]$ is contained in
        $\mathrm{SU}(2)_- \cong \mathrm{Sp}(1)$.
  \item The metric $g_{\mu\nu}[\Theta]$ satisfies the \textbf{anti-self-dual
        condition} $C^+ = 0$.
  \item If additionally $\nabla^\dagger\nabla\Theta = 0$ (vacuum UBT equation),
        then $R_{\mu\nu} = 0$ (Ricci-flat), and by Penrose's theorem,
        $g_{\mu\nu}[\Theta]$ admits a \textbf{curved twistor space description}.
\end{enumerate}
\end{theorem}

\begin{proof}
Part 1 follows from $J_\mu = \Theta^\dagger\partial_\mu\Theta \in \mathfrak{su}(2)_-$
(the current of $\Theta \in \mathrm{SU}(2)_-$).
Part 2 follows from Part 1 and Lemma~\ref{lem:su2_asd}: the holonomy algebra
is contained in $\mathfrak{su}(2)_-$.
Part 3 follows from Part 2 and Theorem~\ref{thm:asd}: holonomy $\subset
\mathrm{Sp}(1) \cong \mathrm{SU}(2)$ implies ASD.
Part 4: $\nabla^\dagger\nabla\Theta = 0$ implies (via the UBT field equations
and the GR recovery chain in \texttt{canonical/gr\_closure/GR\_chain\_summary.tex})
that $G_{\mu\nu} = 0$, hence $R_{\mu\nu} = 0$.  Combined with $C^+ = 0$,
Penrose's nonlinear graviton theorem applies.
\end{proof}

\begin{tcolorbox}[colback=blue!5!white,colframe=blue!50!black,
                  title={Consequence: UBT and Twistor Theory}]
For the $\mathrm{SU}(2)_-$ sector of UBT:
\begin{itemize}
  \item The ASD condition $C^+ = 0$ is proved non-perturbatively \textbf{[L1]}.
  \item The Penrose nonlinear graviton construction applies.
  \item UBT vacuum solutions in this sector correspond to deformations of
        the complex structure of twistor space $\mathbb{PT}$.
  \item This establishes a direct bridge between UBT and Penrose's curved
        twistor theory.
\end{itemize}
\textbf{Note}: This result applies to the instanton sector of UBT, not to
Schwarzschild-type solutions (Petrov type D, not ASD).
\end{tcolorbox}

% ======================================================================
\section{Connection to Twistor Theory}
\label{sec:twistor}
% ======================================================================

\subsection{Penrose Twistor Construction}

If the UBT metric is ASD (or can be restricted to an ASD sector), then
by Penrose's theorem, the curved twistor space $\mathcal{PT}[\Theta]$ is
a deformation of $\mathbb{PT}$ (flat twistor space) parametrised by the
ASD Weyl curvature.

The flat-twistor bridge between UBT and standard Penrose twistors is
documented in \texttt{THEORY\_COMPARISONS/penrose\_twistor/}.  The
curved extension requires:

\begin{enumerate}
  \item The ASD sector to be identified (this document provides the
        linearised version).
  \item The twistor correspondence map $Z^\alpha \leftrightarrow (x,\theta)$
        to be extended to the full curved UBT background.
  \item The Ward theorem to identify gauge fields on twistor space with
        instanton solutions of the UBT field equation on $M^4$.
\end{enumerate}

\subsection{Status of the Twistor Bridge}

\begin{tcolorbox}[colback=blue!5!white,colframe=blue!75!black,
                  title={Summary: Twistor Bridge Status}]
\begin{enumerate}
  \item \textbf{Flat twistor correspondence}: Documented in
        \texttt{THEORY\_COMPARISONS/penrose\_twistor/}.
  \item \textbf{ASD condition (linearised)}: $C^+ = 0$ for
        $\mathrm{SU}(2)_-$ perturbations — \textbf{Proved [L1]}.
  \item \textbf{ASD condition (non-perturbative)}: \textbf{Open [L2]}.
  \item \textbf{Curved twistor space from UBT}: \textbf{Open [L2]}.
  \item \textbf{Ward theorem instanton/gauge field correspondence}:
        \textbf{Open [L2]}.
\end{enumerate}
\end{tcolorbox}

% ======================================================================
\section{Summary}
\label{sec:summary}
% ======================================================================

\begin{tcolorbox}[colback=green!5!white,colframe=green!75!black,
                  title={Summary}]
\begin{enumerate}
  \item The biquaternion algebra $\mathcal{B} = \mathbb{C}\otimes\mathbb{H}$
        contains a natural $\mathrm{SU}(2)_-$ (right-quaternion) subgroup
        that is structurally connected to the ASD condition.

  \item At the linearised level, UBT metric perturbations in the
        $\mathrm{SU}(2)_-$ sector satisfy the ASD condition $C^+ = 0$.
        \textbf{Status: Proved [L1]}.

  \item The full non-perturbative ASD condition requires further analysis
        and is \textbf{Open [L2]}.

  \item The UBT self-dual point $R_\psi = R_t$ (T-duality) is a distinct
        concept from the Weyl-curvature ASD condition; they must not be
        conflated.

  \item A curved twistor description of UBT, via the Penrose nonlinear
        graviton, is structurally plausible given the biquaternion
        $\mathrm{SU}(2)_\pm$ structure, but establishing it rigorously
        requires completing the non-perturbative ASD analysis.
\end{enumerate}

\textbf{Cross-reference}:
\begin{itemize}
  \item \texttt{THEORY\_COMPARISONS/penrose\_twistor/} — flat twistor bridge.
  \item \texttt{research\_tracks/research/schwarzschild\_from\_theta.tex} —
        Schwarzschild $\Theta_0$.
  \item \texttt{canonical/gr\_closure/GR\_chain\_summary.tex} — GR recovery status.
  \item DERIVATION\_INDEX.md --- Gravitational Physics section.
\end{itemize}
\end{tcolorbox}

\end{document}
